\chapter{2018年北京卷（理）}

\section{单选题}

\begin{problem}
已知集合 \(A=\{x\mid |x|<2\}\)，\(B=\{-2,0,1,2\}\)，则 \(A\cap B=\)（\quad）.

\choices
    {\(\{0,1\}\)}
    {\(\{-1,0,1\}\)}
    {\(\{-2,0,1,2\}\)}
    {\(\{-1,0,1,2\}\)}
\end{problem}

\begin{answer}
A
\end{answer}

\begin{solution}
由 \( |x|<2 \) 得 \(-2<x<2\). 在集合 \(B=\{-2,0,1,2\}\) 中，只有 \(0\) 和 \(1\) 满足这个不等式，所以
\(A\cap B=\{0,1\}\)，故选 A.
\end{solution}

\begin{problem}
在复平面内，复数 \(\dfrac{1}{1-\i}\) 的共轭复数对应的点位于（\quad）.

\choices
    {第一象限}
    {第二象限}
    {第三象限}
    {第四象限}
\end{problem}

\begin{answer}
D
\end{answer}

\begin{solution}
\(\dfrac{1}{1-\i}=\dfrac{1+\i}{(1-\i)(1+\i)}=\dfrac{1+\i}{2}\).
因此其共轭复数为 \(\dfrac{1-\i}{2}\)，对应点的坐标为 \(\left(\dfrac12,-\dfrac12\right)\)，横坐标为正、纵坐标为负，故位于第四象限，选 D.
\end{solution}

\begin{problem}
执行如图所示的程序框图，输出的 \(s\) 值为（\quad）.

\texfigure[max width=0.32\linewidth]{img_repaint/2018/beijing_science/q3_fig1.tex}

\choices
    {\(\dfrac12\)}
    {\(\dfrac56\)}
    {\(\dfrac76\)}
    {\(\dfrac{7}{12}\)}
\end{problem}

\begin{answer}
B
\end{answer}

\begin{solution}
程序开始时 \(k=1\)，\(s=1\). 第一次循环后，\(s=1-\dfrac12=\dfrac12\)，且 \(k=2\)；此时 \(k\geq3\) 不成立，继续循环. 第二次循环后，\(s=\dfrac12+\dfrac13=\dfrac56\)，且 \(k=3\)；此时 \(k\geq3\) 成立，输出 \(s=\dfrac56\)，故选 B.
\end{solution}

\begin{problem}
“十二平均律”是通用的音律体系，明代朱载堉最早用数学方法计算出半音比例，为这个理论的发展做出了重要贡献. 十二平均律将一个纯八度音程分成十二份，依次得到十三个单音，从第二个单音起，每一个单音的频率与它的前一个单音的频率的比都等于 \(2^{\frac1{12}}\). 若第一个单音的频率为 \(f\)，则第八个单音的频率为（\quad）.

\choices
    {\(\sqrt[3]{2}\,f\)}
    {\(\sqrt[3]{2^2}\,f\)}
    {\(\sqrt[12]{2^5}\,f\)}
    {\(\sqrt[12]{2^7}\,f\)}
\end{problem}

\begin{answer}
D
\end{answer}

\begin{solution}
从第一个单音到第八个单音共经过 \(7\) 次相邻频率比，因此第八个单音的频率为
\(f\left(2^{\frac1{12}}\right)^7=2^{\frac7{12}}f=\sqrt[12]{2^7}\,f\).
故选 D.
\end{solution}

\begin{problem}
某四棱锥的三视图如图所示，在此四棱锥的侧面中，直角三角形的个数为（\quad）.

\texfigure[max width=6.0cm]{img_repaint/2018/beijing_science/q5_fig1.tex}

\choices
    {\(1\)}
    {\(2\)}
    {\(3\)}
    {\(4\)}
\end{problem}

\begin{answer}
C
\end{answer}

\begin{solution}
由三视图还原四棱锥为 \(P-ABCD\)，可取坐标
\(P(0,2,2)\)，\(A(0,2,0)\)，\(B(0,0,0)\)，\(C(1,0,0)\)，\(D(2,2,0)\).
四个侧面是 \(\triangle PAB\)、\(\triangle PBC\)、\(\triangle PCD\) 和 \(\triangle PDA\). 有
\(\overrightarrow{PA}=(0,0,-2)\)，\(\overrightarrow{AB}=(0,-2,0)\)，\(\overrightarrow{AD}=(2,0,0)\). 因为
\(\overrightarrow{PA}\cdot\overrightarrow{AB}=0\)，\(\overrightarrow{PA}\cdot\overrightarrow{AD}=0\)，所以 \(\triangle PAB\) 和 \(\triangle PDA\) 都是直角三角形.

又 \(\overrightarrow{BP}=(0,2,2)\)，\(\overrightarrow{BC}=(1,0,0)\)，且 \(\overrightarrow{BP}\cdot\overrightarrow{BC}=0\)，所以 \(\triangle PBC\) 是直角三角形. 对于 \(\triangle PCD\)，三个顶点处的相邻边向量内积分别为
\(\overrightarrow{PC}\cdot\overrightarrow{PD}=6\)，\(\overrightarrow{CP}\cdot\overrightarrow{CD}=3\)，\(\overrightarrow{DP}\cdot\overrightarrow{DC}=2\)，均不为 \(0\)，所以 \(\triangle PCD\) 不是直角三角形. 因此直角三角形共有 \(3\) 个，选 C.
\end{solution}

\begin{problem}
设 \(\bs a\)，\(\bs b\) 均为单位向量，则“\(\left\lvert \bs a-3\bs b\right\rvert=\left\lvert 3\bs a+\bs b\right\rvert\)”是“\(\bs a\perp \bs b\)”的（\quad）.

\choices
    {充分而不必要条件}
    {必要而不充分条件}
    {充分必要条件}
    {既不充分也不必要条件}
\end{problem}

\begin{answer}
C
\end{answer}

\begin{solution}
因为 \(\bs a\)，\(\bs b\) 均为单位向量，所以
\(\left\lvert\bs a-3\bs b\right\rvert^2=1+9-6\bs a\cdot\bs b\)，并且
\(\left\lvert3\bs a+\bs b\right\rvert^2=9+1+6\bs a\cdot\bs b\).
两长度相等当且仅当两平方相等，即 \(\bs a\cdot\bs b=0\). 对单位向量，\(\bs a\cdot\bs b=0\) 等价于 \(\bs a\perp\bs b\)，所以该条件是“\(\bs a\perp\bs b\)”的充分必要条件，选 C.
\end{solution}

\begin{problem}
在平面直角坐标系中，记 \(d\) 为点 \(P(\cos\theta,\sin\theta)\) 到直线 \(x-my-2=0\) 的距离，当 \(\theta\)，\(m\) 变化时，\(d\) 的最大值为（\quad）.

\choices
    {\(1\)}
    {\(2\)}
    {\(3\)}
    {\(4\)}
\end{problem}

\begin{answer}
C
\end{answer}

\begin{solution}
令 \(r=\sqrt{1+m^2}\)，则 \(r\geq1\). 点 \(P(\cos\theta,\sin\theta)\) 到直线的距离为
\(d=\dfrac{\left\lvert\cos\theta-m\sin\theta-2\right\rvert}{\sqrt{1+m^2}}\).
由柯西不等式，\(\left\lvert\cos\theta-m\sin\theta\right\rvert\leq r\)，从而
\[
d\leq\frac{r+2}{r}=1+\frac2r\leq3
\]
当 \(m=0\)，\(\theta=\pi\) 时，\(P=(-1,0)\)，直线为 \(x-2=0\)，距离为 \(3\)，所以最大值确为 \(3\)，选 C.
\end{solution}

\begin{problem}
设集合 \(A=\{(x,y)\mid x-y\ge1,ax+y>4,x-ay\le2\}\)，则（\quad）.

\choices
    {对任意实数 \(a\)，\((2,1)\in A\)}
    {对任意实数 \(a\)，\((2,1)\notin A\)}
    {当且仅当 \(a<0\) 时，\((2,1)\notin A\)}
    {当且仅当 \(a\le \dfrac32\) 时，\((2,1)\notin A\)}
\end{problem}

\begin{answer}
D
\end{answer}

\begin{solution}
将 \((x,y)=(2,1)\) 代入三个不等式，分别得到
\(2-1\geq1\)，\(2a+1>4\)，\(2-a\leq2\).
前一个恒成立，后两个分别等价于 \(a>\dfrac32\) 和 \(a\geq0\). 因此
\((2,1)\in A\) 当且仅当 \(a>\dfrac32\)，所以 \((2,1)\notin A\) 当且仅当 \(a\leq\dfrac32\)，选 D.
\end{solution}

\section{填空题}

\begin{problem}
设 \(\{a_n\}\) 是等差数列，且 \(a_1=3\)，\(a_2+a_5=36\)，则 \(\{a_n\}\) 的通项公式为\fillinblank{}.
\end{problem}

\begin{answer}
\(a_n=6n-3\)
\end{answer}

\begin{solution}
设等差数列的公差为 \(d\)，则 \(a_2=3+d\)，\(a_5=3+4d\). 由题意 \((3+d)+(3+4d)=36\)，
解得 \(d=6\). 因而
\(a_n=a_1+(n-1)d=3+6(n-1)=6n-3\).
\end{solution}

\begin{problem}
在极坐标系中，直线 \(\rho\cos\theta+\rho\sin\theta=a\)（\(a>0\)）与圆 \(\rho=2\cos\theta\) 相切，则 \(a=\fillinblank{}\).
\end{problem}

\begin{answer}
\(1+\sqrt2\)
\end{answer}

\begin{solution}
令 \(x=\rho\cos\theta\)，\(y=\rho\sin\theta\). 直线方程化为 \(x+y=a\). 圆的方程满足
\(\rho^2=2\rho\cos\theta=2x\)，即 \(x^2+y^2=2x\)，所以圆心为 \((1,0)\)，半径为 \(1\).

直线与圆相切，当且仅当圆心到直线的距离等于半径，即
\(\dfrac{|1-a|}{\sqrt2}=1\).
故 \(a=1\pm\sqrt2\). 由 \(a>0\)，舍去 \(1-\sqrt2\)，得 \(a=1+\sqrt2\).
\end{solution}

\begin{problem}
设函数 \(f(x)=\cos\!\left(\omega x-\dfrac{\pi}{6}\right)\)（\(\omega>0\)），若 \(f(x)\le f\!\left(\dfrac{\pi}{4}\right)\) 对任意实数 \(x\) 都成立，则 \(\omega\) 的最小值为\fillinblank{}.
\end{problem}

\begin{answer}
\(\dfrac23\)
\end{answer}

\begin{solution}
因为对任意实数 \(x\) 都有 \(f(x)\leq f\left(\dfrac\pi4\right)\)，而余弦函数的最大值为 \(1\)，所以必须有
\(f\left(\frac\pi4\right)=1\).
这等价于存在整数 \(k\)，使得
\(\dfrac{\omega\pi}{4}-\dfrac\pi6=2k\pi\).
于是 \(\omega=\dfrac23+8k\). 由 \(\omega>0\)，有 \(k\geq0\)，故最小值在 \(k=0\) 时取得，为 \(\dfrac23\). 此时 \(f\left(\dfrac\pi4\right)=1\)，条件确实成立.
\end{solution}

\begin{problem}
若 \(x\)，\(y\) 满足 \(x+1\le y\le 2x\)，则 \(2y-x\) 的最小值是\fillinblank{}.
\end{problem}

\begin{answer}
3
\end{answer}

\begin{solution}
由存在 \(y\) 满足 \(x+1\leq y\leq2x\)，得 \(x+1\leq2x\)，所以 \(x\geq1\). 对固定的 \(x\)，表达式 \(2y-x\) 随 \(y\) 增大而增大，故取 \(y=x+1\) 时最小. 此时
\[
2y-x=2(x+1)-x=x+2\geq3
\]
当 \(x=1\)，\(y=2\) 时取到等号，所以最小值为 \(3\).
\end{solution}

\begin{problem}
能说明“若 \(f(x)>f(0)\) 对任意的 \(x\in(0,2]\) 都成立，则 \(f(x)\) 在 \([0,2]\) 上是增函数”为假命题的一个函数是\fillinblank{}.
\end{problem}

\begin{answer}
\(f(x)=\sin x\)
\end{answer}

\begin{solution}
取 \(f(x)=\sin x\). 因为 \(0<x\leq2<\pi\)，所以 \(f(x)=\sin x>0=\sin0=f(0)\)，题设条件成立. 但 \(\dfrac\pi2<2<\pi\)，且
\(f\left(\dfrac\pi2\right)=1>\sin2=f(2)\)，这与在 \([0,2]\) 上为增函数矛盾. 因此该函数可以说明命题为假.
\end{solution}

\begin{problem}
已知椭圆 \(M:\dfrac{x^2}{a^2}+\dfrac{y^2}{b^2}=1\)（\(a>b>0\)），双曲线 \(N:\dfrac{x^2}{m^2}-\dfrac{y^2}{n^2}=1\). 若双曲线 \(N\) 的两条渐近线与椭圆 \(M\) 的四个交点及椭圆 \(M\) 的两个焦点恰为一个正六边形的顶点，则椭圆 \(M\) 的离心率为\fillinblank{}；双曲线 \(N\) 的离心率为\fillinblank{}.
\end{problem}

\begin{answer}
\(\sqrt3-1\)；\(2\)
\end{answer}

\begin{solution}
设椭圆的焦点为 \(F_1(-c,0)\)，\(F_2(c,0)\)，则 \(c^2=a^2-b^2\)，椭圆离心率为 \(e=\dfrac ca\). 六个点关于原点中心对称，且两个焦点在横轴上，因此正六边形的另外四个顶点分别为
\(\left(\dfrac c2,\dfrac{\sqrt3c}{2}\right)\)，\(\left(\dfrac c2,-\dfrac{\sqrt3c}{2}\right)\)，\(\left(-\dfrac c2,\dfrac{\sqrt3c}{2}\right)\)，\(\left(-\dfrac c2,-\dfrac{\sqrt3c}{2}\right)\).
所以双曲线的两条渐近线斜率的绝对值为 \(\sqrt3\)，即 \(\dfrac nm=\sqrt3\). 取其中一个交点代入椭圆方程，得
\[
\frac{c^2}{4a^2}+\frac{3c^2}{4b^2}=1
\]
又 \(b^2=a^2-c^2=a^2(1-e^2)\)，故
\[
\frac{e^2}{4}+\frac{3e^2}{4(1-e^2)}=1
\]
整理得 \(e^4-8e^2+4=0\). 于是
\(e^2=4\pm2\sqrt3\). 因为 \(0<e<1\)，所以
\(e^2=4-2\sqrt3=(\sqrt3-1)^2\)，从而 \(e=\sqrt3-1\).

双曲线的离心率为
\[
e_N=\sqrt{1+\frac{n^2}{m^2}}=\sqrt{1+3}=2
\]
\end{solution}

\section{解答题}

\begin{problem}
在 \(\triangle ABC\) 中，\(a=7\)，\(b=8\)，\(\cos B=-\dfrac17\).

\begin{enumerate}
    \item 求 \(\angle A\)；
    \item 求 \(AC\) 边上的高.
\end{enumerate}
\end{problem}

\begin{solution}
\begin{enumerate}
    \item 因为 \(B\in(0,\pi)\)，所以
    \[
    \sin B=\sqrt{1-\cos^2B}=\sqrt{1-\dfrac1{49}}=\dfrac{4\sqrt3}{7}
    \]
    由正弦定理，
    \[
    \sin A=\dfrac{a\sin B}{b}=\dfrac{7\cdot\dfrac{4\sqrt3}{7}}8=\dfrac{\sqrt3}{2}
    \]
    又因为 \(a<b\)，所以 \(A<B\)，而 \(\cos B<0\) 表明 \(B\) 是钝角，从而 \(A\) 是锐角. 因此 \(A=\dfrac\pi3\).
    \item 设 \(c=AB\). 由余弦定理，\(b^2=a^2+c^2-2ac\cos B\)，即 \(64=49+c^2+2c\)，所以 \(c^2+2c-15=0\)，即 \((c-3)(c+5)=0\). 因为 \(c>0\)，故 \(c=3\).

    设 \(B\) 到 \(AC\) 的高为 \(h\)，则 \(h=c\sin A=3\sin\dfrac\pi3=\dfrac{3\sqrt3}{2}\).
\end{enumerate}
\end{solution}

\begin{problem}
如图，在三棱柱 \(ABC-A_1B_1C_1\) 中，\(CC_1\perp\) 平面 \(ABC\)，\(D\)，\(E\)，\(F\)，\(G\) 分别为 \(AA_1\)，\(AC\)，\(A_1C_1\)，\(BB_1\) 的中点，\(AB=BC=\sqrt5\)，\(AC=AA_1=2\).

\begin{enumerate}
    \item 求证：\(AC\perp\) 平面 \(BEF\)；
    \item 求二面角 \(B-CD-C_1\) 的余弦值；
    \item 证明：直线 \(FG\) 与平面 \(BCD\) 相交.
\end{enumerate}

\texfigure[max width=0.62\linewidth]{img_repaint/2018/beijing_science/q16_fig1.tex}
\end{problem}

\begin{solution}
由 \(E\) 是 \(AC\) 的中点，得 \(AE=EC=1\). 又因为 \(AB=BC\)，所以在等腰三角形 \(ABC\) 中，\(BE\perp AC\)，且
    \(BE=\sqrt{AB^2-AE^2}=\sqrt{5-1}=2\).
由于三棱柱的侧棱互相平行，且 \(CC_1\perp\) 平面 \(ABC\)，所以 \(EF\parallel CC_1\)，\(EF\perp\) 平面 \(ABC\)，并且 \(EF=CC_1=AA_1=2\). 以 \(E\) 为原点，以 \(EB\)，\(EC\)，\(EF\) 所在直线为坐标轴建立空间直角坐标系，则
    \(B(2,0,0)\)，\(C(0,1,0)\)，\(D(0,-1,1)\)，\(F(0,0,2)\)，\(C_1(0,1,2)\)，\(G(2,0,1)\).

\begin{enumerate}
    \item 因为 \(BE\perp AC\)，且 \(EF\perp\) 平面 \(ABC\)，所以 \(EF\perp AC\). 直线 \(BE\) 与 \(EF\) 相交于 \(E\)，且 \(BE\)，\(EF\) 都在平面 \(BEF\) 内，故 \(AC\perp\) 平面 \(BEF\).
    \item 在平面 \(BCD\) 内，取垂直于棱 \(CD\) 且指向点 \(B\) 的向量
    \[
    \bs u=\overrightarrow{CB}-\dfrac{\overrightarrow{CB}\cdot\overrightarrow{CD}}{|\overrightarrow{CD}|^2}\overrightarrow{CD}
    \]
    由 \(\overrightarrow{CB}=(2,-1,0)\)，\(\overrightarrow{CD}=(0,-2,1)\)，得 \(\bs u=(2,-1,0)-\dfrac25(0,-2,1)=\left(2,-\dfrac15,-\dfrac25\right)\).
    在平面 \(CDC_1\) 内，取垂直于棱 \(CD\) 且指向点 \(C_1\) 的向量
    \[
    \bs v=\overrightarrow{CC_1}-\dfrac{\overrightarrow{CC_1}\cdot\overrightarrow{CD}}{|\overrightarrow{CD}|^2}\overrightarrow{CD}
    =(0,0,2)-\dfrac25(0,-2,1)=\left(0,\dfrac45,\dfrac85\right)
    \]
    因而二面角 \(B-CD-C_1\) 的余弦值为
    \[
    \cos\angle(\bs u,\bs v)=\dfrac{\bs u\cdot\bs v}{|\bs u||\bs v|}
    =\dfrac{-\dfrac45}{\sqrt{\dfrac{21}{5}\cdot\dfrac{16}{5}}}=-\dfrac1{\sqrt{21}}
    \]
    \item 平面 \(BCD\) 的一个法向量为 \(\bs n=(1,2,4)\)，因为 \(\bs n\cdot\overrightarrow{CB}=0\)，\(\bs n\cdot\overrightarrow{CD}=0\). 该平面的方程为 \(x+2y+4z=2\). 直线 \(FG\) 上的点可表示为 \(X=F+t\overrightarrow{FG}=(2t,0,2-t)\).
    将其代入平面 \(BCD\) 的方程，得 \(2t+4(2-t)=2\)，解得 \(t=3\). 因此直线 \(FG\) 与平面 \(BCD\) 相交，交点为 \((6,0,-1)\).
\end{enumerate}
\end{solution}

\begin{problem}
电影公司随机收集了电影的有关数据，经分类整理得到下表：

\begin{center}
\begin{tabular}{c|cccccc}
\hline
电影类型 & 第一类 & 第二类 & 第三类 & 第四类 & 第五类 & 第六类 \\
\hline
电影部数 & 140 & 50 & 300 & 200 & 800 & 510 \\
好评率 & 0.4 & 0.2 & 0.15 & 0.25 & 0.2 & 0.1 \\
\hline
\end{tabular}
\end{center}

好评率是指：一类电影中获得好评的部数与该类电影的部数的比值.

假设所有电影是否获得好评相互独立.

\begin{enumerate}
    \item 从电影公司收集的电影中随机选取 \(1\) 部，求这部电影是获得好评的第四类电影的概率；
    \item 从第四类电影和第五类电影中各随机选取 \(1\) 部，估计恰有 \(1\) 部获得好评的概率；
    \item 假设每类电影得到人们喜欢的概率与表格中该类电影的好评率相等，用“\(\xi_k=1\)”表示第 \(k\) 类电影得到人们喜欢，“\(\xi_k=0\)”表示第 \(k\) 类电影没有得到人们喜欢（\(k=1\)，\(2\)，\(3\)，\(4\)，\(5\)，\(6\)）. 写出方差 \(D(\xi_1)\)，\(D(\xi_2)\)，\(D(\xi_3)\)，\(D(\xi_4)\)，\(D(\xi_5)\)，\(D(\xi_6)\) 的大小关系.
\end{enumerate}
\end{problem}

\begin{solution}
\begin{enumerate}
    \item 电影总数为 \(140+50+300+200+800+510=2000\). 第四类电影中获得好评的部数估计为 \(200\times0.25=50\)，所以从全部电影中随机选取 \(1\) 部，该电影是获得好评的第四类电影的概率估计为 \(\dfrac{50}{2000}=0.025\).
    \item 第四类电影获得好评的概率为 \(0.25\)，第五类电影获得好评的概率为 \(0.2\). 由独立性，恰有 \(1\) 部获得好评的概率估计为 \(0.25(1-0.2)+(1-0.25)0.2=0.35\).
    \item 设第 \(k\) 类电影得到人们喜欢的概率为 \(p_k\)，则 \(\xi_k\) 服从两点分布，且 \(D(\xi_k)=p_k(1-p_k)\). 依次计算得 \(D(\xi_1)=0.24\)，\(D(\xi_2)=0.16\)，\(D(\xi_3)=0.1275\)，\(D(\xi_4)=0.1875\)，\(D(\xi_5)=0.16\)，\(D(\xi_6)=0.09\).
    因此
    \[
    D(\xi_6)<D(\xi_3)<D(\xi_2)=D(\xi_5)<D(\xi_4)<D(\xi_1)
    \]
\end{enumerate}
\end{solution}

\begin{problem}
设函数 \(f(x)=\bigl[ax^2-(4a+1)x+4a+3\bigr]\e^x\).

\begin{enumerate}
    \item 若曲线 \(y=f(x)\) 在点 \((1,f(1))\) 处的切线与 \(x\) 轴平行，求 \(a\)；
    \item 若 \(f(x)\) 在 \(x=2\) 处取得极小值，求 \(a\) 的取值范围.
\end{enumerate}
\end{problem}

\begin{solution}
\begin{enumerate}
    \item 对函数求导，得
    \[
    f'(x)=\bigl[2ax-(4a+1)+ax^2-(4a+1)x+4a+3\bigr]\e^x
    =\bigl[ax^2-(2a+1)x+2\bigr]\e^x
    \]
    因为切线与 \(x\) 轴平行，所以 \(f'(1)=(1-a)\e=0\)，因而 \(a=1\).
    \item 将导函数因式分解，得 \(f'(x)=(x-2)(ax-1)\e^x\).
    因为 \(\e^x>0\)，只需考察 \((x-2)(ax-1)\) 在 \(x=2\) 两侧的符号.

    当 \(a>\dfrac12\) 时，\(2a-1>0\)，所以在 \(x=2\) 的左侧，\(f'(x)<0\)，在右侧，\(f'(x)>0\)，从而 \(f(x)\) 在 \(x=2\) 处取得极小值.

    当 \(a<\dfrac12\) 时，\(2a-1<0\)，两侧导数符号相反，左侧为正、右侧为负，故 \(x=2\) 处为极大值点.

    当 \(a=\dfrac12\) 时，\(f'(x)=\dfrac12(x-2)^2\e^x\ge0\)，函数在 \(x=2\) 附近单调增加，因而不在 \(x=2\) 处取得极小值.

    综上，\(a\) 的取值范围为 \(\left(\dfrac12,+\infty\right)\).
\end{enumerate}
\end{solution}

\begin{problem}
已知抛物线 \(C:y^2=2px\) 经过点 \(P(1,2)\). 过点 \(Q(0,1)\) 的直线 \(l\) 与抛物线 \(C\) 有两个不同的交点 \(A\)，\(B\)，且直线 \(PA\) 交 \(y\) 轴于 \(M\)，直线 \(PB\) 交 \(y\) 轴于 \(N\).

\begin{enumerate}
    \item 求直线 \(l\) 的斜率的取值范围；
    \item 设 \(O\) 为原点，\(\overrightarrow{QM}=\lambda\overrightarrow{QO}\)，\(\overrightarrow{QN}=\mu\overrightarrow{QO}\)，求证：\(\dfrac1\lambda+\dfrac1\mu\) 为定值.
\end{enumerate}
\end{problem}

\begin{solution}
\begin{enumerate}
    \item 竖直直线 \(x=0\) 与抛物线只有一个交点，不能满足题设的两个不同交点条件. 因此设直线 \(l\) 的斜率为 \(k\)，则其方程为 \(y=kx+1\). 因为抛物线经过 \(P(1,2)\)，所以 \(4=2p\)，即 \(p=2\)，抛物线方程为 \(y^2=4x\).

    联立直线与抛物线方程，得 \(k^2x^2+(2k-4)x+1=0\). 要有两个不同的交点，必须有 \(k\ne0\)，且判别式满足 \(\Delta=(2k-4)^2-4k^2=16(1-k)>0\). 因此由两交点条件得到 \(k<1\) 且 \(k\ne0\).
    还需保证题目中的点 \(M\)，\(N\) 为有限交点. 当 \(k=-3\) 时，\(l\) 的方程为 \(y=-3x+1\)，与抛物线联立得 \(9x^2-10x+1=0\)，故两个交点中有 \((1,-2)\). 由于 \(P=(1,2)\)，过 \(P\) 与该交点的直线为 \(x=1\)，不与 \(y\) 轴相交，所以此时 \(M\) 或 \(N\) 不存在. 反之，当 \(k<1\)，\(k\ne0\)，\(-3\) 时，两个交点不同，且交点不在直线 \(x=1\) 上，故 \(PA\)，\(PB\) 均与 \(y\) 轴有有限交点. 因此斜率范围为 \((-\infty,-3)\cup(-3,0)\cup(0,1)\).
    \item 由上面的斜率范围，\(k\ne1\)，\(-3\)，所以 \(y_i\ne2\)，\(-2\)，上述各分式均有意义. 设 \(A(x_1,y_1)\)，\(B(x_2,y_2)\). 由 \(A\)，\(B\) 在抛物线上，得 \(x_i=\dfrac{y_i^2}{4}\). 于是 \(\dfrac{y_i-2}{x_i-1}=\dfrac{y_i-2}{\dfrac{y_i^2}{4}-1}=\dfrac4{y_i+2}\). 直线 \(PA\) 与 \(y\) 轴交于 \(M(0,y_M)\)，故 \(y_M-2=\dfrac4{y_1+2}(0-1)\)，\(y_M=\dfrac{2y_1}{y_1+2}\). 同理，\(y_N=\dfrac{2y_2}{y_2+2}\). 由 \(\overrightarrow{QM}=\lambda\overrightarrow{QO}\) 和 \(\overrightarrow{QN}=\mu\overrightarrow{QO}\)，得 \(\lambda=1-y_M=\dfrac{2-y_1}{y_1+2}\)，\(\mu=1-y_N=\dfrac{2-y_2}{y_2+2}\).

    由于 \(y_i=kx_i+1=\dfrac{k y_i^2}{4}+1\)，所以 \(y_1\)，\(y_2\) 是方程 \(ky^2-4y+4=0\) 的两个根，从而 \(y_1+y_2=\dfrac4k\)，\(y_1y_2=\dfrac4k\).
    因此
    \[
    \dfrac1\lambda+\dfrac1\mu
    =\dfrac{y_1+2}{2-y_1}+\dfrac{y_2+2}{2-y_2}
    =\dfrac{8-2y_1y_2}{4-2(y_1+y_2)+y_1y_2}
    =\dfrac{8-\dfrac8k}{4-\dfrac8k+\dfrac4k}=2
    \]
    所以 \(\dfrac1\lambda+\dfrac1\mu\) 为定值 \(2\).
\end{enumerate}
\end{solution}

\begin{problem}
设 \(n\) 为正整数，集合 \(A=\left\{(t_1,t_2,\ldots,t_n)\mid t_i\in\{0,1\},\ i=1,2,\ldots,n\right\}\).
对于集合 \(A\) 中的任意元素 \(\alpha=(x_1,x_2,\ldots,x_n)\) 和 \(\beta=(y_1,y_2,\ldots,y_n)\)，记 \(M(\alpha,\beta)=\frac12\sum_{i=1}^n\bigl(x_i+y_i-|x_i-y_i|\bigr)\).

\begin{enumerate}
    \item 当 \(n=3\) 时，若 \(\alpha=(1,1,0)\)，\(\beta=(0,1,1)\)，求 \(M(\alpha,\alpha)\) 和 \(M(\alpha,\beta)\) 的值；
    \item 当 \(n=4\) 时，设 \(B\) 是 \(A\) 的子集，且满足：对于 \(B\) 中的任意元素 \(\alpha\)，\(\beta\)，当 \(\alpha\)，\(\beta\) 相同时，\(M(\alpha,\beta)\) 是奇数；当 \(\alpha\)，\(\beta\) 不同时，\(M(\alpha,\beta)\) 是偶数. 求集合 \(B\) 中元素个数的最大值；
    \item 给定不小于 \(2\) 的 \(n\)，设 \(B\) 是 \(A\) 的子集，且满足：对于 \(B\) 中的任意两个不同的元素 \(\alpha\)，\(\beta\)，\(M(\alpha,\beta)=0\). 写出一个集合 \(B\)，使其元素个数最多，并说明理由.
\end{enumerate}
\end{problem}

\begin{solution}
注意到当 \(x_i\)，\(y_i\in\{0,1\}\) 时，
\[
\dfrac12\left(x_i+y_i-|x_i-y_i|\right)=
\begin{cases}
1,&x_i=y_i=1,\\
0,&\text{其余情形}.
\end{cases}
\]
所以 \(M(\alpha,\beta)\) 等于 \(\alpha\)，\(\beta\) 在相同位置上同时取 \(1\) 的个数.

\begin{enumerate}
    \item 向量 \(\alpha=(1,1,0)\) 中有两个位置取值为 \(1\)，所以 \(M(\alpha,\alpha)=2\). 向量 \(\alpha\)，\(\beta\) 只有第二个位置同时取值为 \(1\)，所以 \(M(\alpha,\beta)=1\).
    \item 当 \(\alpha=\beta\) 时，\(M(\alpha,\alpha)\) 是 \(\alpha\) 中 \(1\) 的个数. 因而 \(B\) 中的元素只能有奇数个 \(1\)，即只能是含有一个 \(1\) 或含有三个 \(1\) 的向量.

    记含有一个 \(1\) 且该位置为第 \(i\) 个的向量为 \(\alpha_i\)，记只有第 \(i\) 个位置为 \(0\) 的向量为 \(\beta_i\)，其中 \(i=1\)，\(2\)，\(3\)，\(4\). 对于不同的下标，\(M(\alpha_i,\alpha_j)=0\)，\(M(\beta_i,\beta_j)=2\).
    而
    \[
    M(\alpha_i,\beta_j)=
    \begin{cases}
    0,&i=j,\\
    1,&i\ne j.
    \end{cases}
    \]
    因此，若 \(B\) 同时含有某个 \(\alpha_i\) 和某个 \(\beta_j\)，则必须有 \(i=j\). 若 \(B\) 同时含有两类向量，则两类中都至多含有一个元素，故此时 \(|B|\le2\). 若只含有一类向量，则 \(|B|\le4\).

    取
    \[
    B=\left\{(1,0,0,0),(0,1,0,0),(0,0,1,0),(0,0,0,1)\right\}
    \]
    任意两个不同元素的 \(M\) 值都为 \(0\)，满足题意. 所以集合 \(B\) 中元素个数的最大值为 \(4\).
    \item 取
    \[
    B=\left\{(0,0,\ldots,0),(1,0,\ldots,0),(0,1,0,\ldots,0),\ldots,(0,0,\ldots,0,1)\right\}
    \]
    其中任意两个不同元素不可能在同一位置同时取 \(1\)，故任意两个不同元素 \(\alpha\)，\(\beta\) 都满足 \(M(\alpha,\beta)=0\)，且 \(|B|=n+1\).

    另一方面，若 \(B\) 中有 \(r\) 个非零元素，则由任意两个不同元素的 \(M\) 值为 \(0\)，这些非零元素中取值为 \(1\) 的位置两两不重复. 每个非零元素至少占用一个位置，所以 \(r\le n\). 零元素至多有一个，因而 \(|B|\le n+1\). 所以上述集合的元素个数最多，为 \(n+1\).
\end{enumerate}
\end{solution}
